% Общий преамбул для конспекта. Собирается через XeLaTeX/LuaLaTeX.
% Адаптировано из личного конспекта Analysis I
% (/Users/vyorkin/projects/learn/lean4/analysis/notes/preamble.tex) — см. там же
% STYLE.md с полным описанием всех макросов и типографских решений.

\usepackage{polyglossia}
\setmainlanguage{russian}
\setotherlanguage{english}

\usepackage{fontspec}
\setmainfont{PT Serif}
% Без Scale=MatchLowercase Menlo визуально заметно крупнее и жирнее PT Serif
% при том же кегле — инлайн-код "прыгает" на фоне обычного текста.
% MatchLowercase подгоняет размер моноширинного шрифта по x-height основного.
\setmonofont{Menlo}[Scale=MatchLowercase]
\newfontfamily\cyrillicfont{PT Serif}
\newfontfamily\cyrillicfonttt{Menlo}[Scale=MatchLowercase]

% По умолчанию LaTeX увеличивает пробел после точки, которую считает концом
% предложения (англоязычная традиция) — из-за этого один и тот же на вид
% пробел после разных предложений мог визуально отличаться по ширине,
% особенно когда предложение заканчивается на символ из математического
% режима (как "$\Omega$."). В русской типографике пробел после точки везде
% одинаковый — ровно это и делает \frenchspacing.
\frenchspacing

\usepackage{amsmath, amssymb, amsthm, mathtools}

% amsthm по умолчанию продолжает текст доказательства на той же строке, что
% и заголовок "Доказательство." — здесь после заголовка всегда перенос
% строки. Копия \newenvironment{proof} из amsthm.sty с добавленным \newline.
\makeatletter
\renewenvironment{proof}[1][\proofname]{\par
  \pushQED{\qed}%
  \normalfont \topsep6\p@\@plus6\p@\relax
  \trivlist
  \item[\hskip\labelsep
        \itshape
    #1\@addpunct{.}]\hspace{0pt}\\\ignorespaces
}{%
  \popQED\endtrivlist\@endpefalse
}
\makeatother
\usepackage{microtype}
\usepackage{enumitem}
\usepackage{graphicx}
\graphicspath{{images/}}
\usepackage{xcolor}
\usepackage{framed} % рамка \leftbar для окружения intuition, см. ниже
\usepackage{tikz} % минималистичные иллюстрации к отдельным доказательствам
\usetikzlibrary{calc} % арифметика над координатами ($(x)+(dx,dy)$)

% Документ читается с экрана, не печатается — поля уже, чем в книжной
% вёрстке, и одинаковые с обеих сторон (документ не переплетается).
\usepackage[margin=2.7cm]{geometry}

% Чуть больше межстрочного интервала, чем "книжный" одинарный — легче
% воспринимать длинные доказательства и оставлять место для мысленных пометок.
\linespread{1.08}

% Приглушённый серый — для служебных пометок (ссылка на Lean-исходник),
% чтобы они не конкурировали по контрасту с математическим текстом.
\definecolor{leangray}{gray}{0.58}

% Приглушённый синий — для кликабельных ссылок на другие определения/теоремы
% (\bookref). hyperref подключён с [hidelinks], то есть сам он ссылки никак
% не выделяет — цвет здесь единственный сигнал, что по тексту можно
% кликнуть, поэтому красим текст ссылки вручную сами.
\definecolor{bookblue}{RGB}{37,84,150}

% Тот же приглушённый серый, что и leangray — для линии слева и подписи
% окружения intuition (см. ниже). Отдельное имя цвета, а не переиспользование
% leangray напрямую, чтобы два визуально разных по смыслу элемента (ссылка на
% Lean-исходник и блок с интуицией) не были жёстко связаны одним именем —
% если в будущем понадобится развести их цвета, менять придётся только здесь.
\definecolor{intuitiongray}{gray}{0.55}

% Компактные, но не слипшиеся списки по умолчанию — пригодится для шагов
% доказательства/условий упражнений.
\setlist{itemsep=0.2em, topsep=0.5em}

\usepackage[hidelinks]{hyperref}
\usepackage{xurl} % разрешает перенос \path/\url по любому символу, не только после "_"/"."
\usepackage{cleveref}

% У Menlo нет глифа "≫" (U+226B, композиция морфизмов в стиле mathlib) — внутри
% \texttt/\code этот символ иначе печатается как пустой прямоугольник ("Missing
% character"). newunicodechar подменяет именно этот символ на математический \gg
% из основного шрифта формул, независимо от того, в каком текущем шрифте он
% встретился (в частности, внутри моноширинных Lean-сигнатур в \leansource/\code).
\usepackage{newunicodechar}
\newunicodechar{≫}{\ensuremath{\gg}}

% hyperref по умолчанию делает \url/\path курсивным, если курсивен окружающий
% текст (\urlstyle{same}) — для моноширинного шрифта это выглядит плохо
% (синтетический наклон). Код должен быть прямым всегда.
\urlstyle{tt}

% Заголовки разделов вида "1.2" — как в книге
\renewcommand\thesection{\thechapter.\arabic{section}}

% Внутренние (немые) amsthm-окружения. Их автосчётчик не используется по
% значению — он лишь механизм amsthm для печати номера. Сам номер мы каждый
% раз подставляем вручную через \renewcommand\theaxiomEnv, см. ниже.
%
% Везде используется свой стиль "book": жирный заголовок, ПРЯМОЙ текст тела
% (стандартный "plain" даёт курсивное тело, и курсив на длинных абзацах и
% особенно в моноширинном шрифте читается плохо), а формулировка начинается
% с новой строки (последний параметр \newtheoremstyle — "отступ после
% заголовка" — здесь буквально \newline, а не длина). Без этого заголовок
% и первые слова формулировки делят одну строку с оставшейся её частью,
% и при полной выключке текста пробел в этой строке иногда растягивается
% заметно сильнее, чем на соседних строках, — с новой строки эта строка
% всегда обычная строка бегущего текста, без короткого жирного довеска.
\newtheoremstyle{book}{3pt}{3pt}{}{}{\bfseries}{.}{\newline}{}
\theoremstyle{book}
\newtheorem{axiomEnv}{Аксиома}
\newtheorem{theoremEnv}[axiomEnv]{Теорема}
\newtheorem{propositionEnv}[axiomEnv]{Утверждение}
\newtheorem{lemmaEnv}[axiomEnv]{Лемма}
\newtheorem{corollaryEnv}[axiomEnv]{Следствие}
\newtheorem{definitionEnv}[axiomEnv]{Определение}
\newtheorem{exampleEnv}[axiomEnv]{Пример}
\newtheorem{exerciseEnv}[axiomEnv]{Упражнение}
\newtheorem{assumptionEnv}[axiomEnv]{Предположение}
\newtheorem{remarkEnv}[axiomEnv]{Замечание}
\newtheorem*{axiomEnv*}{Аксиома}
\newtheorem*{theoremEnv*}{Теорема}
\newtheorem*{propositionEnv*}{Утверждение}
\newtheorem*{lemmaEnv*}{Лемма}
\newtheorem*{corollaryEnv*}{Следствие}
\newtheorem*{definitionEnv*}{Определение}
\newtheorem*{exampleEnv*}{Пример}
\newtheorem*{exerciseEnv*}{Упражнение}
\newtheorem*{assumptionEnv*}{Предположение}
\newtheorem*{remark}{Замечание}

% Дословная нумерация книги. Использование:
%   \begin{definition}{1.1.1}{категория} ... \end{definition}
% Первый аргумент — номер БУКВАЛЬНО как в комментарии над Lean-объявлением
% (см. CLAUDE.md, «Нумерация в Lean-файлах»), второй — заголовок в скобках
% (можно оставить {}). В этой книге определения/теоремы/леммы/примеры
% нумеруются арабскими цифрами по разделу (Definition 1.1.1), а упражнения —
% римскими (Exercise 1.4.vi) — единой сквозной автонумерации по разделу тут
% не годится ни для одного из двух случаев, поэтому номер всегда указывается
% вручную, ровно в том виде, в каком он стоит в Lean-комментарии (включая
% римские цифры для упражнений).
% Каждое окружение заодно ставит \label{префикс:номер} — это даёт возможность
% ссылаться на него настоящей кликабельной ссылкой через \bookref (см.
% ниже), без ручной возни с метками.
%
% \nopagebreak в конце запрещает разрыв страницы сразу после формулировки —
% без него между "Определение 1.1.1." и следующим за ним "Доказательство."
% может влезть граница страницы, что рвёт нить рассуждения.
\newenvironment{axiom}[2]{\renewcommand\theaxiomEnv{#1}\begin{axiomEnv}[#2]\label{ax:#1}}{\end{axiomEnv}\nopagebreak}
\newenvironment{theorem}[2]{\renewcommand\theaxiomEnv{#1}\begin{theoremEnv}[#2]\label{thm:#1}}{\end{theoremEnv}\nopagebreak}
\newenvironment{proposition}[2]{\renewcommand\theaxiomEnv{#1}\begin{propositionEnv}[#2]\label{prop:#1}}{\end{propositionEnv}\nopagebreak}
\newenvironment{lemma}[2]{\renewcommand\theaxiomEnv{#1}\begin{lemmaEnv}[#2]\label{lem:#1}}{\end{lemmaEnv}\nopagebreak}
\newenvironment{corollary}[2]{\renewcommand\theaxiomEnv{#1}\begin{corollaryEnv}[#2]\label{cor:#1}}{\end{corollaryEnv}\nopagebreak}
\newenvironment{definition}[2]{\renewcommand\theaxiomEnv{#1}\begin{definitionEnv}[#2]\label{def:#1}}{\end{definitionEnv}\nopagebreak}
\newenvironment{example}[2]{\renewcommand\theaxiomEnv{#1}\begin{exampleEnv}[#2]\label{ex:#1}}{\end{exampleEnv}\nopagebreak}
\newenvironment{exercise}[2]{\renewcommand\theaxiomEnv{#1}\begin{exerciseEnv}[#2]\label{exr:#1}}{\end{exerciseEnv}\nopagebreak}
\newenvironment{assumption}[2]{\renewcommand\theaxiomEnv{#1}\begin{assumptionEnv}[#2]\label{assum:#1}}{\end{assumptionEnv}\nopagebreak}

% Нумерованное замечание (в отличие от беззвёздного \remark выше, который печатается
% вовсе без номера) — для тех случаев, когда само замечание пронумеровано в книге.
% Название оставлено отдельным от \remark, чтобы не менять смысл уже использующихся
% в разделах \begin{remark}...\end{remark}.
\newenvironment{remarknum}[2]{\renewcommand\theaxiomEnv{#1}\begin{remarkEnv}[#2]\label{rem:#1}}{\end{remarkEnv}\nopagebreak}

% Беззвёздная версия — для утверждений, которые в книге не пронумерованы
% (например, приведены как рассуждение в прозе, без формальной метки).
% Аргумент — заголовок в скобках, номер не печатается.
\newenvironment{axiom*}[1]{\begin{axiomEnv*}[#1]}{\end{axiomEnv*}\nopagebreak}
\newenvironment{theorem*}[1]{\begin{theoremEnv*}[#1]}{\end{theoremEnv*}\nopagebreak}
\newenvironment{proposition*}[1]{\begin{propositionEnv*}[#1]}{\end{propositionEnv*}\nopagebreak}
\newenvironment{lemma*}[1]{\begin{lemmaEnv*}[#1]}{\end{lemmaEnv*}\nopagebreak}
\newenvironment{corollary*}[1]{\begin{corollaryEnv*}[#1]}{\end{corollaryEnv*}\nopagebreak}
\newenvironment{definition*}[1]{\begin{definitionEnv*}[#1]}{\end{definitionEnv*}\nopagebreak}
\newenvironment{example*}[1]{\begin{exampleEnv*}[#1]}{\end{exampleEnv*}\nopagebreak}
\newenvironment{exercise*}[1]{\begin{exerciseEnv*}[#1]}{\end{exerciseEnv*}\nopagebreak}
\newenvironment{assumption*}[1]{\begin{assumptionEnv*}[#1]}{\end{assumptionEnv*}\nopagebreak}

% Ссылка на конкретную лемму/теорему в Lean-файле того же раздела.
% \path (из url.sty, с xurl) даёт моноширинный текст, переносимый по любому
% символу — длинные идентификаторы не вылезают за поля, и не нужно
% экранировать подчёркивания вручную.
\newcommand{\leanref}[1]{\path{#1}}

% Короткая вставка кода прямо в тексте (не длинный идентификатор — для тех
% есть \leanref). \mbox запрещает перенос строки внутри — короткий фрагмент
% вроде "zero : Nat" не должен разрываться пополам на границе строки.
\newcommand{\code}[1]{\mbox{\texttt{#1}}}

% Внутри инлайн-формулы запрещает перенос строки ВНУТРИ данной скобочной
% группы — например, в цепочке "$x \iff \nosplit{(x \in A \lor x \in B)}
% \iff y$", чтобы перенос по ширине не разрезал "(x ∈ A ∨ x ∈ / B)" пополам.
% Не \mbox: \mbox замораживает ширину группы совершенно жёстко, и это лишает
% microtype/обычную выключку возможности чуть сжать-растянуть соседние
% пробелы — из-за этого короткая скобка неожиданно вызывала Overfull hbox
% (текст на пару мм вылезал за поле), хотя визуально казалось, что перенос
% просто ушёл на следующую строку. \relpenalty/\binoppenalty — пенальти
% именно за перенос на отношениях (\in, \subseteq, класс Rel) и бинарных
% связках (\land, \lor, класс Bin) — TeX по умолчанию считает их разрешёнными
% точками переноса ВНУТРИ формулы. Подняв оба до 10000 локально (в пределах
% группы {...}), запрещаем перенос именно там, но не мешаем обычной выключке
% искать другую точку (например, перед всей группой, на внешнем \iff).
\newcommand{\nosplit}[1]{{\relpenalty=10000\binoppenalty=10000 #1}}

% Вынесенная строка с именем Lean-теоремы/леммы вместе с её формальной
% сигнатурой — ставится последней внутри окружения, после
% доказательства/формулировки. #1 — идентификатор, #2 — явные параметры и
% формулировка (без неявных {..}-параметров), в родном синтаксисе Lean.
% Флеш-лефт, а не hfill: сигнатура обычно длиннее одного идентификатора и
% переносится на несколько строк, вправо это смотрелось бы плохо.
%
% Две вещи для читаемости длинных сигнатур:
%   1) \raggedright — это код, а не проза. Полная выключка растягивала
%      межсловные пробелы на перенесённых строках так же уродливо, как и
%      в обычном тексте, только здесь это особенно бросается в глаза;
%   2) \hangindent — если сигнатура не помещается и переносится, продолжение
%      уходит с отступом, а не жмётся к левому краю вровень с "Lean:".
% Перенос происходит только когда сигнатура реально не помещается — короткие
% остаются на одной строке с именем.
% Запрещает разрыв страницы между строками ЭТОГО абзаца — без него длинная
% сигнатура, перенесённая на 2-3 строки через \leanbreak, могла разъехаться
% по границе страницы. \par\nobreak в начале уже запрещает разрыв ПЕРЕД этим
% \leansource — см. \nopagebreak в определениях axiom/theorem/... выше.
%
% Обёрнуто в \parbox, а не просто \clubpenalty/\widowpenalty/\interlinepenalty
% — те применяются только к разрывам ВНУТРИ одного абзаца, а \leanbreak
% внутри \raggedright не создаёт второй "строки" того же абзаца: \raggedright
% переопределяет \\ на \@centercr, которая сама вызывает \par — то есть
% \leanbreak на самом деле начинает НОВЫЙ абзац. \parbox делает всю сигнатуру
% одним неделимым боксом: TeX либо целиком помещает её на текущую страницу,
% либо целиком переносит на следующую.
\newcommand{\leansource}[2]{\par\nobreak\smallskip\noindent
  \parbox[t]{\linewidth}{\raggedright\hangindent=2em\hangafter=1
  \footnotesize\color{leangray}Lean:~{\leanref{#1}}\ \texttt{#2}}\par}

% Ручной перенос строки внутри сигнатуры \leansource{}{...} — для длинных
% сигнатур, которые не помещаются в одну строку. Ставится в точке
% семантического разрыва (после имени теоремы, если аргументов нет, или
% перед финальным ":", отделяющим гипотезы от заключения) — не там, где
% попало у автоматического переноса по ширине строки.
\newcommand{\leanbreak}{\\\hspace*{2em}}

% Единообразный разбор случаев внутри \begin{proof}...\end{proof}:
%   \caselabel{$\Omega \in \Omega$} Тогда ...
\newcommand{\caselabel}[1]{\medskip\noindent\textbf{Случай #1.}\\}

% То же самое, но с произвольной подписью — для шагов индукции и похожих
% структурных частей доказательства:
%   \prooflabel{База} $n = 0$: ...
%   \prooflabel{Переход} ...
\newcommand{\prooflabel}[1]{\medskip\noindent\textbf{#1.}\\}

% Кликабельная ссылка на другое определение/теорему/... по её метке (см.
% префиксы выше — ax:, thm:, prop:, lem:, cor:, def:, ex:, exr:, assum:).
% Падеж и формулировку текста задаём сами вторым аргументом, макрос не
% гадает.
%   По \bookref{def:1.1.1}{определению 1.1.1}, ...
\newcommand{\bookref}[2]{\hyperref[#1]{\color{bookblue}#2}}

% Обоснование шага доказательства — мелкой серой пометкой сразу после
% содержательного утверждения, вместо ведущего "По определению X, ...". Само
% рассуждение остаётся обычным текстом на переднем плане, ссылка на
% определение/лемму отходит на второй план (как и \leansource):
%   $f$ — изоморфизм.\by{\bookref{def:1.1.9}{определение 1.1.9}}
\newcommand{\by}[1]{\ {\footnotesize\color{leangray}(#1)}}

% Блок с интуитивным объяснением — тонкая вертикальная линия слева (как
% цитата), мелкая серая подпись "Интуиция" сверху, дальше обычный текст без
% заливки и без изменения цвета. Ставится после формулировки теоремы/
% определения (пояснить, что она значит и зачем нужна в такой форме) или
% перед/внутри \begin{proof} (пояснить замысел доказательства, выбор
% конкретного свидетеля и т.п.) — вне самих theorem-подобных окружений, чтобы
% не мешать \leansource оставаться последней строкой внутри них.
%
% framed.sty (\MakeFramed/\FrameCommand), а не ручной \vrule в minipage —
% minipage не умеет переносить содержимое через границу страницы, а
% framed.sty именно для этого и существует.
\newenvironment{intuition}{%
  \def\FrameCommand{{\color{intuitiongray}\vrule width 1pt}\hspace{10pt}}%
  \MakeFramed{\advance\hsize-\width \FrameRestore}%
  \noindent{\footnotesize\color{intuitiongray}\textbf{Интуиция}}\par\smallskip
  \noindent\ignorespaces
}{%
  \par\endMakeFramed\par\medskip
}
